\section{Incremental Deployment Strategy}
\label{sec:deployment}

% GUIDANCE: Show how the constellation can be built incrementally, providing
% value from day one.

% 5.1 Deployment Phases
%   - Ring-by-ring deployment. Each new ring placed to bisect the largest
%     remaining gap (binary-search analogy).
%   - Phase 1: single ring at ~1.3 AU midpoint, ~357 sats, ~2 Mbps.
%     Already matches current Mars relay capability.
%   - Phase 2: second ring, total ~944 sats, ~31 Mbps (15x improvement).
%   - Phases 3--6: each ring roughly doubles throughput.
%     By ring 6: ~2,208 sats, ~197 Mbps.
%   - Further densification (more sats per ring) + more rings -> 1 Gbps.
%   - INCLUDE FIGURE: staircase chart of cumulative throughput vs sats deployed.

% 5.2 Ring Addition Order
%   - Discuss why bisection is optimal: the bottleneck is always the widest
%     gap, so inserting a ring there gives the largest marginal improvement.
%   - Relate to the R_opt formula: as S grows, you naturally want more rings
%     proportional to sqrt(S).

% 5.3 Orbital Mechanics
%   - Heliocentric circular orbits between 1.0 and 1.52 AU.
%   - Delta-v from LEO: 1--6 km/s beyond Earth escape depending on target ring.
%   - Each Starship launch delivers ~20 satellites to a single ring.
%   - In-plane phasing maneuvers to distribute satellites around the ring.
%
% Target length: ~1.5 pages.
